% !TEX program = lualatex
% Normative native TeX/TikZ specimen for learning-dynamics-2406.
% The scientific content is illustrative and original; it does not reproduce
% data or scientific content from arXiv:2406.07856v1.
\documentclass[tikz,border=6pt]{standalone}

\usepackage{learning-dynamics-2406}
\usetikzlibrary{arrows.meta,backgrounds,calc,positioning}

% Scientific parameters -------------------------------------------------------
\def\TauG{1.35}
\def\EvenTarget{0.82}
\def\OddTarget{0.28}

\tikzset{paneltag/.append style={anchor=south west}}

\begin{document}
\begin{tikzpicture}[x=1cm,y=1cm]
  \begin{scope}[on background layer]
    \fill[Canvas] (-0.15,-0.15) rectangle (17.95,7.55);
  \end{scope}

  % Figure and phase ribbons --------------------------------------------------
  \fill[NeutralBar,rounded corners=3pt] (0,6.82) rectangle (17.8,7.42);
  \node[majorheading] at (8.9,7.12) {Adaptive field response};

  \fill[TrainBar,rounded corners=3pt] (0,6.05) rectangle (11.35,6.57);
  \node[sectiontitle] at (5.675,6.31) {Training};
  \fill[RetrieveBar,rounded corners=3pt] (11.72,6.05) rectangle (17.8,6.57);
  \node[sectiontitle] at (14.76,6.31) {Retrieval};

  % Panel (a): imposed motion -------------------------------------------------
  \node[paneltag] at (0.05,5.22) {(a)};
  \node[paneltitle] at (2.65,5.35) {imposed motion};

  \begin{scope}[on background layer]
    \foreach \rx/\ry/\op in {1.65/0.82/0.035,1.30/0.64/0.055,
                                0.98/0.48/0.085,0.68/0.33/0.13,
                                0.38/0.19/0.20}{
      \fill[FieldCore,opacity=\op] (4.05,3.40) ellipse (\rx cm and \ry cm);
    }
  \end{scope}

  \node[entityi] (ai) at (1.35,3.40) {$i$};
  \node[entityj] (aj) at (4.05,3.40) {$j$};
  \draw[trainarrow] ($(ai.north)+(0,0.20)$) -- ++(1.05,0);
  \draw[trainarrow] ($(aj.north)+(0,0.20)$) -- ++(1.05,0);

  \draw[axisline] (1.35,2.55) -- (4.05,2.55);
  \draw[axisline] (1.35,2.46) -- (1.35,2.64);
  \draw[axisline] (4.05,2.46) -- (4.05,2.64);
  \node[mathlabel] at (2.70,2.30) {$a$};
  \node[mathlabel] at (2.70,1.43)
    {$\tau_g\dot g=g_{\mathrm{target}}-g$};

  % Panel (b): learned response ----------------------------------------------
  \node[paneltag] at (5.88,5.22) {(b)};
  \node[paneltitle,anchor=west] at (6.60,5.35) {learn interaction};

  \begin{scope}[shift={(6.42,1.22)},x=0.63cm,y=3.45cm]
    \draw[axisline] (0,0) -- (6.15,0);
    \draw[axisline] (0,0) -- (0,0.92);
    \draw[axisline] (0,0) -- (6,0);

    \foreach \x/\lab in {0/0,3/3,6/6}{
      \draw[axisline] (\x,0) -- ++(0,-0.025);
      \node[ticklabel,anchor=north] at (\x,-0.055) {\lab};
    }
    \foreach \y/\lab in {0/0,0.4/0.4,0.8/0.8}{
      \draw[axisline] (0,\y) -- ++(-0.04,0);
      \node[ticklabel,anchor=east] at (-0.10,\y) {\lab};
    }

    \draw[positive]
      plot[domain=0:6,samples=100,smooth]
      (\x,{\EvenTarget*(1-exp(-\x/\TauG))});
    \draw[negative]
      plot[domain=0:6,samples=100,smooth]
      (\x,{\OddTarget*(1-exp(-\x/\TauG))});

    \node[mathlabel,text=InteractionBlue,anchor=south east]
      at (5.9,0.82) {$g^{+}$};
    \node[mathlabel,text=InteractionRed,anchor=north east]
      at (5.9,0.25) {$g^{-}$};
    \node[mathlabel] at (3,-0.20) {$t$};
    \node[annotation,rotate=90] at (-1.05,0.44) {learned response $g$};
  \end{scope}

  % Panel (c): retrieved drift ------------------------------------------------
  \node[paneltag] at (11.88,5.22) {(c)};
  \node[paneltitle,anchor=west] at (12.66,5.35) {adaptive drift};

  \begin{scope}[on background layer]
    \clip (12.45,1.02) rectangle (17.48,5.05);
    \foreach \rx/\ry/\op in {2.25/1.40/0.035,1.80/1.12/0.055,
                                1.35/0.84/0.085,0.92/0.57/0.13,
                                0.48/0.30/0.20}{
      \fill[FieldCore,opacity=\op] (16.18,3.20) ellipse (\rx cm and \ry cm);
    }
  \end{scope}
  \draw[snapshot] (12.45,1.02) rectangle (17.48,5.05);

  \draw[Retrieve,line width=0.65pt,opacity=0.35]
    (13.12,1.62) .. controls (13.75,2.05) and (14.20,2.63) .. (14.63,2.86);
  \draw[retrievearrow]
    (14.60,2.85) .. controls (15.05,3.08) and (15.35,3.18) .. (15.78,3.20);
  \foreach \p/\op in {{(13.12,1.62)}/0.28,{(13.65,2.00)}/0.38,
                         {(14.16,2.55)}/0.52,{(14.65,2.87)}/0.68,
                         {(15.20,3.10)}/0.82}{
    \fill[Retrieve,opacity=\op] \p circle (1.05pt);
  }

  \node[entityi,opacity=0.38] at (13.12,1.62) {$i$};
  \node[entityi] at (15.82,3.20) {$i$};
  \node[entityj] at (16.55,3.20) {$j$};
  \node[annotation,text=Muted,anchor=north west] at (13.24,1.50) {start};
  \node[annotation,anchor=south east] at (17.31,1.18)
    {illustrative; not to scale};
  \node[mathlabel] at (14.97,0.66) {$x$};
  \node[mathlabel,rotate=90] at (12.03,3.04) {$y$};
\end{tikzpicture}
\end{document}
